\section{The character contribution}

The first line of \eqref{eq:nullity-ledger} is a commutative torsion-point
count.  The following calculation identifies it over the ground field and
also fixes the sign in \eqref{eq:D-definition}.

\begin{proposition}[Cyclotomic intersection]\label{prop:character-count}
The number of singular character blocks is
$\Dcycl(\alpha,\beta,\gamma)$.
\end{proposition}

\begin{proof}
For $u,v\in\mu_\ell(k)$, singularity means
\begin{equation}\label{eq:character-equation}
 \alpha+\beta u+\gamma v=0.
\end{equation}
Because $\gamma\ne0$, a choice of $u$ determines the only possible value
\[
 v=-\frac{\alpha+\beta u}{\gamma}.
\]
Since $\ell$ is odd, the condition $v^\ell=1$ becomes
\[
 -\frac{(\alpha+\beta u)^\ell}{\gamma^\ell}=1,
 \qquad\text{or equivalently}\qquad
 (\alpha+\beta u)^\ell+\gamma^\ell=0.
\]
Thus singular characters are in bijection with the common roots of the two
polynomials in \eqref{eq:D-definition}.  The polynomial $t^\ell-1$ is
separable in characteristic $p$ because $p\ne\ell$.  Hence the degree of the
gcd counts those common roots once each.
\end{proof}

\begin{remark}[Arithmetic content]\label{rem:arithmetic-content}
The term $\Dcycl$ is an intersection of the $\ell$-torsion torus with the
affine curve $\alpha+\beta u+\gamma v=0$.  It satisfies
$0\le\Dcycl\le\ell$ and can vary with both primes and the coefficient point.
This is the part of the answer analogous to lattice Fourier analysis in
positive characteristic \citep{Zaidenberg2008}.  The next section accounts
for the nonabelian mass absent from that calculation.
For $(\alpha,\beta,\gamma)=(1,1,1)$, the corresponding resultant is the
classical Wendt determinant \citep{FordJha1993}; we use the gcd only as a
block count and claim no new resultant theory.
\end{remark}
